% SI — PhysBench
\documentclass[11pt,a4paper]{article}
\usepackage[margin=2.4cm]{geometry}
\usepackage{booktabs}
\usepackage{amsmath,amssymb}
\usepackage{hyperref}
\usepackage{microtype}
\usepackage{enumitem}
\setlength{\parskip}{0.4em}
\linespread{1.12}

\title{\bfseries Supplementary Information\\
PhysBench: a verifiable multi-domain benchmark for physics-informed learning}
\author{Sehaj Randhir Singh \\ \small Correspondence: sehajrsinghs@gmail.com}
\date{}
\begin{document}
\maketitle

\section*{SI 1. Parameter ranges}
\label{si:ranges}

Table~\ref{tab:ranges} lists the full parameter ranges and the deterministic
seeding protocol for each domain.

\begin{table}[h]
\centering
\caption{PhysBench parameter ranges.}
\label{tab:ranges}
\small
\begin{tabular}{@{}ll@{}}
\toprule
Domain & Parameters (range) \\
\midrule
beam & $L\in[1,6]$, $P\in[10^3,5\times10^4]$, $E\in[69,210]$\,GPa, $I\in[10^{-7},2\times10^{-5}]$, $h\in[0.05,0.4]$ \\
cantilever & identical ranges to beam (identical token stream) \\
projectile & $v_0\in[5,30]$\,m/s, $\alpha\in[15^\circ,75^\circ]$ \\
pendulum & $L\in[0.5,3]$, $\theta_0\in[5^\circ,60^\circ]$ \\
spring & $k\in[1,50]$, $m\in[0.1,5]$, $A\in[0.2,2]$ \\
rc & $R\in[100,10^4]$, $C\in[10^{-6},10^{-3}]$, $V_0\in[1,12]$ \\
burgers & $\nu\in[0.02,0.1]$, $A\in[0.8,1.8]$, $\sigma\in[0.25,0.45]$ \\
heat2d & mode numbers $k,l\in\{1,2,3\}$, amplitudes $\in[0.5,2]$ \\
\bottomrule
\end{tabular}
\end{table}

\section*{SI 2. Solver and verifier independence}
\label{si:verifiers}

For the trajectory domains, the generator is the exact closed-form solution
(\texttt{physx/sim.py}) and the verifier is a finite-difference
discretization of the governing equation (\texttt{physx/residuals.py}): e.g.\
for the beam, the generator evaluates the closed-form deflection
$w(x) = Px(4x^2-3L^2)/(48EI)$ (piecewise), while the verifier checks the
fourth-derivative equilibrium $EI\,w'''' = -q$ numerically on the predicted
curve. The two code paths share no implementation. For Burgers, the generator
is a high-resolution upwind scheme; the verifier evaluates
$u_t + uu_x - \nu u_{xx}$ by central differences on the predicted grid. For
heat2d, the generator sums the separable-mode solution; the verifier evaluates
$\nabla^2 u - f$ on the predicted grid.

\section*{SI 3. The baseline matrix, per-seed}
\label{si:matrix}

The full 75-run matrix (per-seed training rel-MAE) is committed under
\texttt{physx/models/matrix/}. Table~\ref{tab:seeds} summarizes means and
standard deviations per cell.

\begin{table}[h]
\centering
\caption{Baseline matrix (mean $\pm$ std over 5 seeds).}
\label{tab:seeds}
\small
\begin{tabular}{@{}lccc@{}}
\toprule
Domain & phys & no-physics & MLP \\
\midrule
beam & $0.118\pm0.029$ & $0.153\pm0.048$ & $0.133\pm0.037$ \\
cantilever & $0.120\pm0.027$ & $0.150\pm0.030$ & $0.134\pm0.030$ \\
burgers & $0.008\pm0.001$ & $0.010\pm0.002$ & $0.016\pm0.003$ \\
heat2d & $0.014\pm0.007$ & $0.013\pm0.003$ & $0.020\pm0.007$ \\
projectile & $0.100\pm0.031$ & $0.057\pm0.030$ & $0.068\pm0.025$ \\
\bottomrule
\end{tabular}
\end{table}

\section*{SI 4. Reproducibility}
\label{si:repro}

\begin{verbatim}
python3 -c "from physx import dataset; print(len(dataset.generate('beam', n=256, seed=0)))"
python3 physx/run_matrix.py        # rerun the 75-run baseline matrix
python3 -m unittest physx.test_physx   # 49 tests, incl. dataset/verifier tests
\end{verbatim}

Reproducing the matrix takes under an hour on a 20-core CPU; data generation
is instantaneous and deterministic.

\end{document}
